% ScholarWeft LaTeX export — "document" doc type.
%
% This is NOT a full pandoc --template (see book.tex's own note for why) — it
% is included into the document PREAMBLE via pandoc's --include-in-header,
% alongside pandoc's own default LaTeX template (which already renders the
% TOC, citeproc bibliography, and font-package setup — no need to reproduce
% any of that here).
%
% Every line below is ordinary, editable LaTeX preamble content — copy this
% file (e.g. into your own Export Templates folder) and change, add, or
% remove anything: a journal logo (see \swTitleLogo below), 2-column layout
% (see SWEXTRACLASSOPTIONS below), different heading styles, extra packages,
% etc. Python only ever supplies the note's CONTENT (title/subtitle/author/
% date/abstract/body/citations) — never how it looks.
%
% Available substitutions, filled in by DocumentCompiler.py before this file
% reaches pandoc (NOT pandoc $variable$ syntax — see the note below):
%   SWTOKAUTHOR       — the note's author
%   SWTOKSHORTTITLE   — the note's short title (falls back to the full title)
%   SWTOKTITLE        — the note's title (unused by default; available if
%                        you want it somewhere other than the title block,
%                        e.g. in a running header instead of the short title)
% Available hooks (LaTeX commands DocumentCompiler.py invokes when building
% the title block/page — redefine them below to customize; left empty they
% do nothing):
%   \swTitleLogo      — invoked once, centered, just above the title
% Available marker (a plain LaTeX comment DocumentCompiler.py itself reads —
% inert to LaTeX, meaningful only to the Python side):
%   % SWEXTRACLASSOPTIONS: <comma-separated \documentclass options>
%                        — e.g. "twocolumn" or "12pt". "twoside" is always
%                        added too (required by the running-header
%                        convention below) — you don't need to repeat it.
% Available substitution tokens driven by export checkboxes, not user-edited
% (listed for completeness — moving/removing them changes what those
% checkboxes do):
%   SWTOKFIGURECOUNTER   — "Restart footnote and figure numbering per
%                        chapter": \counterwithin{figure}{section} (document/
%                        article's top-level heading is \section, not
%                        \chapter) when on, empty when off. Footnotes are
%                        deliberately NOT part of this: article-class
%                        footnotes are naturally continuous with no native
%                        per-section reset to hook into (unlike book, where
%                        \chapter resets footnotes to plain "1, 2, 3..."
%                        automatically — see book.tex), and continuous is
%                        the desired behaviour here regardless of the
%                        checkbox, so there's nothing to substitute.
%   SWTOKNEWPAGEHEADING  — "Top-level headings start on a new page": \newpage
%                        when on, empty when off. Sits in \titleformat's
%                        before-code slot for \section only (Heading 2 never
%                        forces a page break, matching DOCX/ODT).
%
% SWTOKAUTHOR / SWTOKSHORTTITLE / SWTOKTITLE are replaced with the document's
% real (LaTeX-escaped) values before this file is handed to pandoc — NOT
% pandoc $variable$ syntax, which LaTeX would read as inline math ($..$)
% since --include-in-header content is raw LaTeX, not run through pandoc's
% templating engine.
%
% Title BLOCK, not title page: matches document.docx/document.odt, whose
% title/subtitle/author/date sit as ordinary paragraphs at the top of page 1
% with no page break before the body — DocumentCompiler.py builds this as
% body content (see export_latex's use of is_book to choose title-block vs.
% title-page). "article" gets the same treatment; only "book" gets a
% dedicated title page, matching book.docx's own section-break-separated
% cover page.
%
% Running header: Author (even pages) / Short Title (odd pages), matching
% the DOCX/ODT "document"/"article" templates' own even/odd header
% convention — requires a two-sided layout (set from Python via
% --metadata classoption=twoside). No header on page 1: LaTeX's kernel
% already puts a document's first page on the `plain` page style regardless
% of \pagestyle, so this falls out for free.

% SWEXTRACLASSOPTIONS:

% ── Customize this template ────────────────────────────────────────────
% Copy this file (e.g. into your own Export Templates folder) and edit any
% of the lines below to get your own fonts/margins/link color/spacing.
%
% Noto Sans here is a VARIABLE font: one file with a continuous `wght` axis
% (100–900) and no separate Bold file. LuaLaTeX exposes that axis through
% fontspec's Weight= key, so BoldFeatures={Weight=700} below is the font's
% REAL 700 instance — not a synthetic fake-bold. Weight=400, the axis
% default, is normal text. (fontspec under XeLaTeX can't address variable
% axes, so the \else branch keeps the old synthetic FakeBold, used only when
% a .tex is compiled by hand with xelatex.) Width=/Slant=/RawAxis= work the
% same way if you ever need them.
\ifluatex
  % ── Unicode GLYPH FALLBACK ──────────────────────────────────────────────
  % LuaLaTeX draws only glyphs the current font HAS. Noto Sans lacks
  % many symbols (→ U+2192, ⚙ U+2699, ✓ U+2713 …) and all EMOJI, which then
  % show as tofu boxes. Build a fallback chain from whichever fonts are
  % installed (each guarded by \IfFontExistsTF) and hand it to luaotfload: a
  % glyph the main font lacks is drawn from the first listed font that has
  % it. Ordered Sans-style symbols first (matched to Noto Sans),
  % then monochrome emoji, CJK, and broad catch-alls for other scripts.
  % NOTE: symbols fonts differ in vertical position — Apple Symbols draws →
  % near the BASELINE, STIX Two Math / Noto Sans Math centre it — so a math/
  % symbols font precedes Apple Symbols. Colour emoji fonts (Apple Color
  % Emoji, Segoe UI Emoji) CANNOT be used (no colour-bitmap support);
  % install monochrome "Noto Emoji" (listed below) or use twemojis.
  \def\swfallbacklist{}
  \IfFontExistsTF{Noto Sans Math}{\xdef\swfallbacklist{\swfallbacklist "Noto Sans Math:mode=node;",}}{}
  \IfFontExistsTF{Noto Sans Symbols 2}{\xdef\swfallbacklist{\swfallbacklist "Noto Sans Symbols 2:mode=node;",}}{}
  \IfFontExistsTF{Noto Sans Symbols}{\xdef\swfallbacklist{\swfallbacklist "Noto Sans Symbols:mode=node;",}}{}
  \IfFontExistsTF{Segoe UI Symbol}{\xdef\swfallbacklist{\swfallbacklist "Segoe UI Symbol:mode=node;",}}{}
  \IfFontExistsTF{STIX Two Math}{\xdef\swfallbacklist{\swfallbacklist "STIX Two Math:mode=node;",}}{}
  \IfFontExistsTF{Apple Symbols}{\xdef\swfallbacklist{\swfallbacklist "Apple Symbols:mode=node;",}}{}
  \IfFontExistsTF{DejaVu Sans}{\xdef\swfallbacklist{\swfallbacklist "DejaVu Sans:mode=node;",}}{}
  \IfFontExistsTF{Symbola}{\xdef\swfallbacklist{\swfallbacklist "Symbola:mode=node;",}}{}
  \IfFontExistsTF{Noto Emoji}{\xdef\swfallbacklist{\swfallbacklist "Noto Emoji:mode=node;",}}{}
  \IfFontExistsTF{Noto Sans CJK SC}{\xdef\swfallbacklist{\swfallbacklist "Noto Sans CJK SC:mode=node;",}}{}
  \IfFontExistsTF{PingFang SC}{\xdef\swfallbacklist{\swfallbacklist "PingFang SC:mode=node;",}}{}
  \IfFontExistsTF{Hiragino Sans GB}{\xdef\swfallbacklist{\swfallbacklist "Hiragino Sans GB:mode=node;",}}{}
  \IfFontExistsTF{WenQuanYi Zen Hei}{\xdef\swfallbacklist{\swfallbacklist "WenQuanYi Zen Hei:mode=node;",}}{}
  \IfFontExistsTF{Segoe UI}{\xdef\swfallbacklist{\swfallbacklist "Segoe UI:mode=node;",}}{}
  \IfFontExistsTF{Arial Unicode MS}{\xdef\swfallbacklist{\swfallbacklist "Arial Unicode MS:mode=node;",}}{}
  \directlua{luaotfload.add_fallback("swfallback", {\swfallbacklist})}
  % ────────────────────────────────────────────────────────────────────────
  \setmainfont{Noto Sans}[BoldFont={Noto Sans},BoldFeatures={Weight=700},
    RawFeature={fallback=swfallback}]
\else
  \setmainfont{Noto Sans}[BoldFont={Noto Sans},BoldFeatures={FakeBold=3.5}]
\fi
\usepackage[letterpaper,margin=1.3in]{geometry}  % page size + margins (try a4paper, or a per-side margin like "top=1in,bottom=1in,left=1.25in,right=1in")
\definecolor{swlinkcolor}{HTML}{156082}        % link/citation color, as a hex RRGGBB (no '#') — matches document.docx/.odt's Hyperlink style
% ── Obsidian callout boxes ──────────────────────────────────────────────────
% scripts/sw-callouts.lua wraps standard Obsidian callouts in this environment,
% passing the type's background colour (a 10% tint, as in Obsidian) and border
% colour. Needs tcolorbox; when it is absent the environment degrades to an
% unboxed block so the content still appears. Tune the padding/rule here.
\IfFileExists{tcolorbox.sty}{\usepackage{tcolorbox}\tcbuselibrary{breakable}}{}
\IfFileExists{changepage.sty}{\usepackage{changepage}}{}
\newif\ifswcalloutindent
\ifcsname adjustwidth\endcsname\swcalloutindenttrue\else\swcalloutindentfalse\fi
\ifcsname newtcolorbox\endcsname
  \newtcolorbox{swcalloutinner}[2]{%
    colback=#1, colframe=#2, boxrule=0.4pt, arc=1.5pt,
    left=6pt, right=6pt, top=4pt, bottom=4pt, boxsep=0pt,
    breakable, before skip=0.3cm, after skip=0.3cm,
  }%
  % Indent the whole box 1.5cm from each margin (matching the DOCX/ODT callout
  % styles' left/right indent), with 0.2cm above/below to set it apart.
  \newenvironment{swcalloutbox}[2]{%
    \ifswcalloutindent\begin{adjustwidth}{1.5cm}{1.5cm}\fi
    \begin{swcalloutinner}{#1}{#2}%
  }{%
    \end{swcalloutinner}%
    \ifswcalloutindent\end{adjustwidth}\fi
  }%
\else
  \newenvironment{swcalloutbox}[2]{}{}%
\fi
\usepackage{setspace}\setstretch{1.25}         % paragraph line spacing — a setspace MULTIPLE (1 = single); this is NOT the same number as the DOCX/ODT 1.15 line-spacing multiple (they measure from a different baseline reference), so tune it by eye.
\SetSinglespace{1.25}                          % FOOTNOTE/FLOAT line spacing — setspace forces footnotes and floats to SINGLE spacing regardless of \setstretch, so they look cramped next to the body; set this to match \setstretch (or lower). Affects footnotes & floats only, not body text.
\raggedbottom                                  % let pages end wherever content naturally does, instead of
                                                % stretching heading/list/paragraph spacing to reach an even
                                                % bottom margin on every page (that stretch is what makes the
                                                % same heading's gap look bigger on one page than another with
                                                % no visible cause). Switch to \flushbottom for even bottom
                                                % margins instead, at the cost of that inconsistent stretch.
% ────────────────────────────────────────────────────────────────────────

% Title-block logo hook — empty by default. Redefine to add a journal/
% institution logo above the title, e.g.:
%   \renewcommand{\swTitleLogo}{\includegraphics[width=3cm]{/path/to/logo.png}\\[1em]}
\newcommand{\swTitleLogo}{}

% ── Arabic / non-Latin RTL scripts (wrapper-free) ───────────────────────
% The LaTeX PDF engine is LuaLaTeX (DocumentCompiler.py's
% find_latex_engine), not XeLaTeX, because of the babel setup below.
% XeLaTeX has no paragraph-level bidi of its own: fontspec's
% [Script=Arabic] shapes and joins each contiguous Arabic run correctly, but
% separate runs on one line (e.g. the two hemistichs of a verse, split by
% \hspace) are still placed left-to-right, so the line comes out reversed.
% XeLaTeX's macro-based fixes — \TeXXeTstate+\beginR/\endR, polyglossia, and
% the standalone bidi package — all either break the letter-joining or die
% with "TeX capacity exceeded [main memory size=5000000]" inside hyperref's
% bidi code at the first \footnote on a real document (all confirmed).
%
% babel's `bidi=basic` is different: it runs the Unicode bidi algorithm over
% LuaTeX's node list, so shaping AND ordering are correct with no macro
% capacity limit. `onchar=ids fonts` then makes it AUTOMATIC: any run of
% Arabic-script characters — including Arabic Supplement / Arabic Extended-A
% and Arabic-derived "Ajami" orthographies (Wolof, Hausa, Fulfulde, ...) — is
% detected by its Unicode script id, switched to the font below, and laid
% out RTL, even inside a roman paragraph. So Arabic mixed into ordinary
% prose needs NO wrapper; `\selectlanguage{arabic}` and sw-poetry.lua's
% `{\arabicfont ...}` block are still honoured when present, but are not
% required for inline text.
%
% The Arabic font is Scheherazade New, scaled to 120% (Scale=1.2): Arabic
% set at the Latin font's nominal size looks smaller because its
% x-height/ascender proportions differ, so the customary correction is a
% small boost. Swap the name(s) below for a different Arabic typeface — a
% static family with a real Bold file (like Scheherazade New) is simplest;
% Noto Naskh Arabic is the fallback when it isn't installed.
%
% Chinese is provided too — CJK needs no bidi (it is LTR), so `onchar` here
% is only for automatic FONT selection; add any other script the same way.
% The nested \IfFontExistsTF guards keep a missing font from breaking the
% export (worst case those characters fall back to the main font).
%
% pandoc's own standalone template already loads babel (also with
% bidi=basic) whenever the note carries a `lang` property; only load it
% ourselves when it isn't already loaded, or the differing options are an
% "option clash" error.

% Right indent for Arabic verse lines. babel maps LaTeX's "left" indents to
% the RTL (right) side, so this controls how far the poetry block sits from
% the RIGHT text margin — the verse environment's own default leaves it
% nearly flush, which reads as a formatting error. Raise it for a deeper
% indent. Affects only Arabic poetry; inline Arabic is unaffected.
\newlength{\swArabicVerseIndent}\setlength{\swArabicVerseIndent}{5em}
\newcommand{\swArabicVerseSetup}{\setlength{\leftmargini}{\swArabicVerseIndent}}

\ifluatex
  \makeatletter
  \@ifpackageloaded{babel}{}{\usepackage[bidi=basic,shorthands=off]{babel}}
  \makeatother
  \babelprovide[import,onchar=ids fonts]{arabic}
  \IfFontExistsTF{Scheherazade New}{%
    \babelfont[arabic]{rm}[Scale=1.2]{Scheherazade New}%
  }{\IfFontExistsTF{Noto Naskh Arabic}{\babelfont[arabic]{rm}[Scale=1.2]{Noto Naskh Arabic}}{}}%
  \babelprovide[import,onchar=ids fonts]{chinese}
  \IfFontExistsTF{Noto Serif CJK SC}{%
    \babelfont[chinese]{rm}{Noto Serif CJK SC}%
  }{\IfFontExistsTF{Songti SC}{\babelfont[chinese]{rm}{Songti SC}}{}}%
  \newcommand{\arabicfont}{\selectlanguage{arabic}\swArabicVerseSetup}
\else
  % XeLaTeX fallback (only reached if a .tex is compiled by hand with
  % xelatex): shapes/joins runs but has no paragraph-level bidi, so mixed
  % inline runs come out reversed. LuaLaTeX is strongly preferred.
  \newfontfamily\swArabicFont[Script=Arabic]{Scheherazade New}
  \newcommand{\arabicfont}{\swArabicFont\swArabicVerseSetup}
\fi

% Bullet-list levels alternate bullet/dash instead of LaTeX's own default
% bullet/dash/asterisk/centered-dot cascade.
\usepackage{enumitem}
\setlist[itemize,1]{label=\textbullet}
\setlist[itemize,2]{label=\textendash}
\setlist[itemize,3]{label=\textbullet}
\setlist[itemize,4]{label=\textendash}

% Heading NUMBERING is decided by DocumentCompiler.py (number_headings +
% latexize_headings), which emits every heading as raw LaTeX: a numbered
% heading becomes \section{...}, while every unnumbered heading becomes the
% starred form \section*/\subsection*/... carrying whatever literal number the
% compiler wrote (e.g. "1.1"). DocumentCompiler.py also emits
% \setcounter{secnumdepth}{0} into the body so no heading can be
% double-numbered. There is therefore no titlesec/titletoc configuration here:
% LaTeX's own defaults style the headings.
SWTOKFIGURECOUNTER

\usepackage{fancyhdr}
\pagestyle{fancy}
\fancyhf{}
\fancyhead[LE]{SWTOKAUTHOR}
\fancyhead[RO]{SWTOKSHORTTITLE}
\fancyfoot[C]{\thepage}
\renewcommand{\headrulewidth}{0.4pt}
% Page 1 (the title block) uses LaTeX's "plain" pagestyle — redefine it to
% keep the centered page-number footer but drop the running header, since
% \_latex_title_block sets \thispagestyle{plain} explicitly (article class
% has no automatic page-1 pagestyle override the way \maketitle would give
% it, since we never call \maketitle — see the title-block note above).
\fancypagestyle{plain}{%
  \fancyhf{}
  \fancyfoot[C]{\thepage}
  \renewcommand{\headrulewidth}{0pt}
}
% Not setting \hypersetup{pdftitle=...} here deliberately: --include-in-header
% content lands BEFORE pandoc's own hyperref setup in its default template, so
% \hypersetup isn't defined yet at this point — the PDF's internal
% pdftitle/pdfauthor metadata is the one thing not replicated for LaTeX
% output (a cosmetic gap, not a content one). colorlinks/linkcolor above
% DOES reach pandoc's own hypersetup call, since that reads --metadata
% variables Python passes separately, resolved against swlinkcolor once
% xcolor's \definecolor above has already run.

